面向嵌入式与边缘设备，日志系统需要在可观测性、写入开销与并发一致性之间取得平衡。针对文本日志在高频场景下格式化成本高、空间冗余大以及多进程共享元数据易发生初始化竞争的问题，本文提出并实现 Optbinlog。该系统采用事件格式描述文件、共享元数据与二进制编码路径，将格式解析从热路径迁出，并通过偏移访问保证跨进程映射一致性。

本文围绕``机制分析---理论推导---分层实验验证''展开：首先调研 Linux ftrace、LTTng、Zephyr logging、systemd-journald、NanoLog、RT-Thread ulog 与 OpenHarmony HiLog Lite 等代表系统，并定位其关键优化路径；随后建立 \(N\)（日志条数）与 \(D\)（并发设备/节点数）驱动下的时间、空间与吞吐模型；最后在统一统计口径下完成工程近似层与严格公平层的对照验证。

实验数据由四个阶段构成：工程近似矩阵、严格语义对齐主套件（single、本地多设备 \(D=5,10,20,50\)）、真实多设备硬件 CRC32C 复测（5/10/15/20 节点）以及初始化竞争验证。结果表明：\thesisterm{binary} 相对 \thesisterm{text} 在本轮纳入语义组上均保持显著正收益；在``Optbinlog 对齐 peer 语义、peer 保持原生存储路径''的严格公平口径下，Optbinlog 在 single 场景已可在时间与吞吐上超过 \thesisterm{nanolog/zephyr}，但空间仍落后；在真实多设备场景中，对 \thesisterm{ulog/syslog} 保持稳定综合优势，对 \thesisterm{hilog} 呈现``时间小幅占优、空间占优、吞吐近似持平''的混合形态，对 \thesisterm{nanolog/zephyr} 仍处负区；\thesisterm{ftrace} 则呈现``空间正收益、时间吞吐负收益''。新增空间扫描进一步显示，\thesisterm{ulog/syslog} 在 \(N\approx 10\) 即出现空间反超，\thesisterm{hilog} 在 \(N\approx 50\) 左右出现交叉，而 \thesisterm{nanolog/zephyr} 在 \(N\le 100000\) 仍未出现交叉点。初始化竞争实验稳定满足``单创建者、其余进程/节点复用已初始化共享文件''的正确性约束。理论趋势与实测结果总体一致。
